% sections/intro_sota.tex -- Introduction and state of the art
\section{Introduction and state of the art}
\label{sec:intro}

\subsection{Universality and field-theoretic approaches}
\label{sec:intro-universality}

Continuous phase transitions in systems with short-range interactions and an $N$-component order parameter with
$O(N)$ symmetry fall into the universality classes of the three-dimensional $O(N)$ model: $N=0$ describes polymers
and self-avoiding walks, $N=1$ uniaxial magnets and the liquid--vapour critical point, $N=2$ planar magnets and the
superfluid transition of $^4$He, $N=3$ isotropic magnets~\cite{PelissettoVicari2002,ZinnJustinBook}. Their critical
exponents are a standard benchmark both for field theory and for non-perturbative methods.

Two perturbative field-theoretic approaches have been pushed to high order. In the $\varepsilon$ expansion the
renormalization-group (RG) functions of $\phi^4$ theory are computed by minimal subtraction in $d=4-\varepsilon$ and
the exponents are expanded in powers of $\varepsilon$~\cite{ZinnJustinBook,KleinertBook}. In fixed-dimension
perturbation theory, proposed by Parisi~\cite{Parisi1980}, one works directly in $d=3$, where the theory is
super-renormalizable, renormalizes it by normalization conditions at zero momentum, and obtains the exponents from
the RG functions at the non-trivial zero of the $\beta$ function. In both approaches the series are asymptotic, with
coefficients growing like $k!$ at large order $k$~\cite{Lipatov1977,BrezinLeGuillouZinnJustin1977}, and must be
resummed, typically by a Borel--Leroy transform followed by a conformal mapping built on the known large-order
behaviour~\cite{LeGuillouZinnJustin1977,LeGuillouZinnJustin1980,GuidaZinnJustin1998}. Fixed dimension has the
advantage that Borel summability is established for the correlation functions of the $d=3$ theory~\cite{MagnenSeneor1977}, although not directly for
the RG functions, while for the $\varepsilon$ series it has to be assumed. The precision of the resummed exponents is limited by the number of known terms, by the
numerical uncertainty of the highest ones, and by systematic effects of the resummation; in fixed dimension these
include the expected non-analyticity of $\beta$ at the fixed point, caused by confluent corrections to
scaling~\cite{Nickel1991,PelissettoVicari2002}.

\subsection{The fixed-dimension series: history and status}
\label{sec:intro-fixed}

We consider the Euclidean $O(N)$ $\phi^4$ theory in $d=3$ in the massive zero-momentum scheme, in which the
renormalized mass $m$, the field normalization and the dimensionless coupling $g$ are fixed by conditions on
$\Gamma^{(2)}_R(0)$, $\partial\Gamma^{(2)}_R/\partial p^2(0)$ and $\Gamma^{(4)}_R(0)=m\,g$. We write the series in
$\ubar=(N+8)u/6$ with $u=g/8\pi$, i.e., in the normalized coupling of
Refs.~\cite{BNGM1976,BNM1978,Phi4tools2024} (the $\tilde u$ of Ref.~\cite{Phi4tools2024}). The computation of the
coefficient of $\ubar^{L+1}$ in $\beta$ requires the four-point function $\Gamma^{(4)}$ at $L$ loops; the coefficients of $\ubar^{L}$
in $\eta$ and $\nu$ require the two-point function $\Gamma^{(2)}$, with its $p^2$ derivative or with $\phi^2$
insertions, at $L$ loops.

\paragraph{Nickel and collaborators.}
A numerical evaluation programme was first carried out by Baker, Nickel, Green and Meiron~\cite{BNGM1976}, who took the series for
$N=1$ from three to six loops~\cite{Ribeca2003thesis}, i.e., $\beta$ to $\ubar^{7}$ and $\eta$, $\nu$ to $\ubar^{6}$;
Baker, Nickel and Meiron~\cite{BNM1978} analysed the corresponding series for $N=0$--$3$. The key technical ingredient was that one-loop bubbles, triangles and boxes in $d=3$
have closed forms~\cite{Nickel1978}, so that each graph could be integrated in momentum space with these one-loop
subgraphs replaced by effective vertices. The individual graph values were compiled, together with their weights, in
an unpublished Guelph report~\cite{NMB1977}. The 1042 diagrams needed through six loops~\cite{Ribeca2003thesis} are
given to ten decimals when the integral has dimension at most three and to eight decimals when it has dimension five
or six, with errors ``generally less than $\pm5$ in the last digit quoted'' and up to $\pm20$ in exceptional
cases. The whole evaluation took about
20~hours on an IBM 370/155 and 3~hours on a CDC 6600~\cite{NMB1977}. The integral of each graph was set up by hand, and the
paper on the numerical methods announced in Refs.~\cite{BNGM1976,NMB1977,BNM1978} never appeared. Because the compilation
separates integrals from symmetry factors, it was reused for decades to build the series of other models without
recomputing any integral, for instance for cubic anisotropy~\cite{CarmonaPelissettoVicari2000}; see
Ref.~\cite{PelissettoVicari2002} for a review. Murray and Nickel later extended $\eta$ and $\eta_{\phi^2}$, hence
$\nu$, to seven loops for $N=0$--$3$~\cite{MurrayNickel1991}; this report is also unpublished, its method is
undocumented, and its coefficients are quoted, to about seven significant digits, in the appendix of
Ref.~\cite{GuidaZinnJustin1998}. The resummation of these series by Le~Guillou and Zinn-Justin and by Guida and
Zinn-Justin~\cite{LeGuillouZinnJustin1977,LeGuillouZinnJustin1980,GuidaZinnJustin1998} produced the standard
fixed-dimension estimates of the $O(N)$ exponents.

\paragraph{Automation.}
By 2003 the six-loop integrals had not been reproduced by an independent group for a quarter of a century, and the
next order involves several thousand graphs. References~\cite{Ribeca2003thesis,GuidaRibeca2006} therefore set out to automate
the whole computation: graph generation, weights and $O(N)$ factors, and an automatic search for the cheapest
effective-vertex parametrization of every amplitude, followed by deterministic adaptive cubature. All seven-loop
amplitudes were parametrized with at most nine integration variables~\cite{Ribeca2003thesis}, but the seven-loop
numerical evaluation was not completed at the time.

\paragraph{Phi4tools.}
Sberveglieri and Spada~\cite{Phi4tools2024} (the Phi4tools compilation) recomputed all zero-momentum graphs with up to eight vertices, that is,
$\Gamma^{(4)}$ to seven loops and $\Gamma^{(2)}$ to eight. They essentially followed the strategy proposed in
Ref.~\cite{GuidaRibeca2006}, integrating every graph in momentum space after automatically assigning effective
vertices,
with only slight differences in their own words: separate tools for each step, Mathematica for the integrals of
dimension up to three and the VEGAS algorithm~\cite{Lepage1978} for the rest in place of deterministic adaptive
cubature, some new analytic and numerical effective vertices, and a simpler search for the parametrization of the
amplitudes. Some of the hardest graphs were cross-checked with feyntrop~\cite{feyntrop2024}, an implementation of
tropical Monte Carlo. Their largest relative uncertainties are a few parts in $10^5$ for two- and
four-point graphs and a few parts in $10^4$ for the $p^2$ derivatives of two-point graphs. This provided the first
independent check of Ref.~\cite{NMB1977}: of the 585 six-loop four-point graphs renormalized in the same way, most
values agree to $\pm1$ in the last digit quoted by Nickel et al., with few exceptions, of which only one differs by
about 20~units. From these values
they obtained $\beta$, $\eta$ and $\nu$ to $\ubar^{8}$ for general $N$, and the corresponding series of the cubic
model. For $N=1$ the quoted uncertainties of the $\ubar^8$ coefficients, added linearly over the coefficients of the
$N$ polynomials, are about $4\times10^{-6}$ ($\beta$), $5\times10^{-4}$ ($\eta$) and $5\times10^{-5}$ ($\nu$) in
relative terms. The authors postponed the resummation to future work~\cite{Phi4tools2024}; to our knowledge the
eighth-order fixed-dimension series have not been resummed. Table~\ref{tab:sota} summarizes this history.

\begin{table*}
\caption{\label{tab:sota}%
Perturbative RG functions of the $O(N)$ model: highest known orders. In fixed dimension the term $\ubar^{k}$ of
$\beta$ needs $\Gamma^{(4)}$ at $L=k-1$ loops, and that of $\eta$ and $\nu$ needs $\Gamma^{(2)}$ at $L=k$ loops; the
graph counts are one-particle-irreducible (1PI) graphs without self-loops, one-vertex-reducible ones (with a cut vertex)
included. The precision is the relative
uncertainty of the highest coefficient at $N=1$ or, where no uncertainty is quoted, the number of significant digits
printed (sig.) or the absolute accuracy of the individual integrals (int.). The $\varepsilon$-expansion coefficients are known exactly, as
rational combinations of multiple zeta values and, at seven loops, of one further period that is conjecturally not of
this form~\cite{Schnetz2018}.}
\begin{ruledtabular}
\begin{tabular}{lllll}
Function & Term & Graphs & Reference & Precision \\
\hline
\multicolumn{5}{l}{\emph{Fixed dimension $d=3$, massive zero-momentum scheme}}\\
$\beta$ & $\ubar^{7}$ & $\Gamma^{(4)}$, $L=6$ (660) & \cite{BNGM1976} ($N=1$); \cite{BNM1978,NMB1977} ($N=0$--$3$) & 6 sig.; int.\ $10^{-8}$--$10^{-10}$ \\
        & $\ubar^{8}$ & $\Gamma^{(4)}$, $L=7$ (3986) & \cite{Phi4tools2024}, all $N$ & $4\times10^{-6}$ \\
        & $\ubar^{9}$ & $\Gamma^{(4)}$, $L=8$ (26\,540) & this work, all $N$ & \betaNineRel \\
        & $\ubar^{10}$ & $\Gamma^{(4)}$, $L=9$ (197\,022) & this work, all $N$ & \betaTenRel \\
        & $\ubar^{11}$ & $\Gamma^{(4)}$, $L=10$ (1\,593\,254) & this work, all $N$ & \betaElevenRel \\
$\eta$, $\nu$ & $\ubar^{6}$ & $\Gamma^{(2)}$, $L=6$ (75) & \cite{BNGM1976} ($N=1$); \cite{BNM1978,NMB1977} ($N=0$--$3$) & int.\ $10^{-8}$--$10^{-10}$ \\
        & $\ubar^{7}$ & $\Gamma^{(2)}$, $L=7$ (317) & \cite{MurrayNickel1991,GuidaZinnJustin1998}, $N=0$--$3$ & $\sim7$ sig. \\
        & $\ubar^{8}$ & $\Gamma^{(2)}$, $L=8$ (1622) & \cite{Phi4tools2024}, all $N$ & $5\times10^{-4}$ ($\eta$), $5\times10^{-5}$ ($\nu$) \\
        & $\ubar^{9}$ & $\Gamma^{(2)}$, $L=9$ (9184) & this work, all $N$ & \etaNineRel\ ($\eta$), \nuNineRel\ ($\nu$) \\
        & $\ubar^{10}$ & $\Gamma^{(2)}$, $L=10$ (59\,614) & this work, all $N$ & \etaTenRel\ ($\eta$), \nuTenRel\ ($\nu$) \\
\hline
\multicolumn{5}{l}{\emph{$\varepsilon$ expansion, minimal subtraction: $\beta$, $\gamma_{m^2}$, $\gamma_\phi$}}\\
        & 5 loops & & \cite{ChetyrkinKataevTkachov1981,ChetyrkinGorishnyLarinTkachov1983,GorishnyLarinTkachov1984,Kleinert1991,KleinertBook}, all $N$ & exact \\
        & 6 loops & & \cite{BatkovichChetyrkinKompaniets2016,KompanietsPanzer2017}, all $N$ & exact \\
        & 7 loops & & \cite{Schnetz2018,Schnetz2023}, all $N$ & exact \\
        & 8 loops & $\gamma_\phi$ only & \cite{Schnetz2023}, $N=1$ & exact \\
\end{tabular}
\end{ruledtabular}
\end{table*}

\subsection{The \texorpdfstring{$\varepsilon$}{epsilon} expansion and non-perturbative benchmarks}
\label{sec:intro-eps}

In the $\varepsilon$ expansion the minimal-subtraction RG functions of the $O(N)$ model were extended to five
loops~\cite{ChetyrkinKataevTkachov1981,ChetyrkinGorishnyLarinTkachov1983,GorishnyLarinTkachov1984,Kleinert1991} and
to six~\cite{BatkovichChetyrkinKompaniets2016,KompanietsPanzer2017}, the latter computation also estimating the graphs
without subdivergences to eleven loops and comparing them with the large-order behaviour of $\beta$. Position-space
methods based on graphical functions have since given the $\beta$ function and the anomalous dimensions to seven
loops~\cite{Schnetz2018}, for the $O(N)$ model as well (published for $N=1$, the $O(N)$ results distributed with the
author's software), and for $N=1$ the field anomalous dimension to eight~\cite{Schnetz2023}. Numerical approaches in minimal subtraction remove subdivergences by
subtraction operators acting on the integrands, so that the RG functions follow algebraically from linear sums of
finite integrals~\cite{AdzhemyanKompaniets2011,AdzhemyanKompaniets2014}; the zero-momentum $\Rbar$ operation used here
(Sec.~\ref{sec:methods-rbar}) is the fixed-dimension counterpart of this idea. The six-loop $\varepsilon$ series have been resummed for general
$N$~\cite{KompanietsPanzer2017,KompanietsWiese2020}, and for $N=1$ the seven-loop series and, for $\eta$, the eight-loop
one~\cite{Shapoval2026}. A recent study for $N=1$~\cite{Shapoval2026} finds that the uncertainty of $\eta$ decreases
rapidly with the loop order while those of $\nu$ and of the correction-to-scaling exponent $\omega$ do not, and
reports a slight but systematic tension with the conformal bootstrap.

The most precise determinations of the exponents are now non-perturbative. The numerical conformal bootstrap
determines the Ising exponents with uncertainties far below those of any resummed perturbative series, and
constrains the $O(N)$ exponents as well~\cite{KosPolandSimmonsDuffinVichi2016,PolandRychkovVichi2019}; Monte Carlo
simulations, of the Ising model~\cite{FerrenbergXuLandau2018} and of improved lattice
Hamiltonians~\cite{Hasenbusch2010,Hasenbusch2019,Hasenbusch2020}, also reach high precision. The conformal bootstrap gives $\nu=\rsBenchCBNuNone$ and
$\eta=\rsBenchCBEtaNone$ for $N=1$~\cite{ChangEtAl2025}, $\nu=\rsBenchCBNuNtwo$ and
$\eta=\rsBenchCBEtaNtwo$ for $N=2$~\cite{ChesterEtAl2020}, and $\nu=\rsBenchCBNuNthree$ and $\eta=\rsBenchCBEtaNthree$ for
$N=3$~\cite{ChesterEtAl2021}, in agreement with the Monte Carlo values $\nu=\rsBenchMCNuNtwo$, $\eta=\rsBenchMCEtaNtwo$
for $N=2$~\cite{Hasenbusch2019} and $\nu=\rsBenchMCNuNthree$, $\eta=\rsBenchMCEtaNthree$ for $N=3$~\cite{Hasenbusch2020}.
For $N=1$ in particular, these results make the perturbative series a stringent test of resummation methods, a
test that is only meaningful if the series coefficients and their uncertainties are reliable.

\subsection{This work}
\label{sec:intro-this-work}

Since 1978 each fixed-dimension series has been extended by one order at a
time~\cite{MurrayNickel1991,Phi4tools2024}. Both complete computations~\cite{NMB1977,Phi4tools2024} used one
momentum-space representation per graph, chosen by hand or automatically and integrated to high accuracy, and
obtained the RG functions from the vertex functions of an intermediate scheme through the $Z$ factors (Nickel et
al.\ subtracted, according to Ref.~\cite{Phi4tools2024}, all two-point subgraphs at zero momentum, Phi4tools only the divergent tadpole and sunset
subgraphs). The ninth and tenth orders need 26\,540 and 197\,022 four-point graphs at eight and nine loops, and
9184 and 59\,614 two-point graphs at nine and ten loops, against 3986 and 1622 at the eighth order, each of higher
dimension; the eleventh order of $\beta$ needs 1\,593\,254 four-point graphs at ten loops. Moreover, along this route the order-$L$
contributions largely cancel: at four, five and six loops the intermediate-scheme four-point input to $\beta$ exceeds
the net order-$L$ four-point coefficient $B_L$ (Sec.~\ref{sec:methods}) by factors of about 29, 50 and 74, so that its relative
error must be smaller than the target by the same factor.

In this paper we compute $\beta$, $\eta$ and $\nu$ to $\ubar^{10}$ ($\beta$ to $\ubar^{11}$) for general $N$, and
re-evaluate all lower orders, with a method designed for this scale and for massively parallel hardware: its
precision depends as little as possible on the order, and its cost should grow slowly with the order
(Sec.~\ref{sec:methods-choice}); measured, from seven loops on, the cost grows by a factor of \costgrowthrange\ per loop
(Sec.~\ref{sec:results-performance}).
\begin{enumerate}
\item Through the Callan--Symanzik equation, the coefficients of $\beta$, $\eta$ and $\nu$ follow algebraically from
three functions, $\Gamma^{(4,1)}_R(0)$, $\partial\Gamma^{(2,1)}_R/\partial p^2(0)$ and $\Gamma^{(2,2)}_R(0)$, whose $L$-loop
terms we call $B_L$, $H_L$ and $C_L$ (Sec.~\ref{sec:methods-cs}). Each is a \emph{linear} sum over graphs of finite
integrals, subtracted at zero momentum by Zimmermann's forest
formula~\cite{BogoliubovParasiuk1957,Hepp1966,Zimmermann1969,Zavialov1990}, the
$\Rbar$ operation of Sec.~\ref{sec:methods-rbar}. No $Z$ factors and no series reversion are needed, and the
cancellations that the $Z$-factor route performs between graphs take place inside each integrand. We call this
route A.
\item At zero external momenta and unit masses every graph, with all its forest terms, is a single parametric
integral over a simplex, built from the Kirchhoff polynomials of the graph and of its subgraphs and quotient
graphs~\cite{BergereZuber1974,BergereLam1976,BrownKreimer2013}. Every integrand is sampled by tropical, Hepp-sector Monte
Carlo~\cite{Volkov2017,Borinsky2023tropical,Panzer2022Hepp,feyntrop2024,Borinsky2025global}, with a density adapted to the
forest terms, on any number of CPUs or GPUs from the same source code.
\item The computation is automatic, from graph generation to series coefficients: there is no per-graph choice of
momentum routing, effective vertices or integration variables. Since every graph carries its $O(N)$ polynomial, one
run gives every $N$.
\item The results are checked against exact low-order values, against the published series, against the large-$N$
limit and the large-order behaviour, and, at the ninth and tenth orders, against a second, independent route (route B)
through the vertex functions of an intermediate scheme, its $Z$ factors and the reversion of the coupling series. The
precision of route B, \routeBrelmin--\routeBrelmax\ relative, is far coarser than that of the new coefficients, and the
coefficient of $\ubar^{11}$ in $\beta$ is checked only by the internal statistics, the large-$N$ limit and the
large-order behaviour.
\end{enumerate}
The results are therefore an independent determination of all known orders: the renormalization route, the
integral representation and the numerics share nothing with Refs.~\cite{NMB1977,Phi4tools2024} except the graph
lists, which our own generator reproduces wherever both exist (Sec.~\ref{sec:methods-checks}). The published coefficients combine our new sums with, at each lower order, the more
precise of our sum and the one derived from Ref.~\cite{Phi4tools2024}; with our own lower orders only, they change
by at most \ledgeronlyzmax\ standard deviations (Sec.~\ref{sec:results-new}).

Runs on NVIDIA H200 GPUs reproduce the known per-order sums from
two to seven loops for all three functions within three standard deviations, with error estimates that the
within-run statistics confirm to a few per cent (Sec.~\ref{sec:results-validation}).
At six loops, with the computing time we allocated to it, the method reproduces the published values with similar or
better precision. From seven loops on, where the graphs are many and their integrals of high dimension, it overtakes
the per-graph treatment: all three seven-loop sums are now known more precisely than
from Phi4tools, by factors of about \precBseven, \precHseven\ and \precCseven\ for $B_7$, $H_7$ and $C_7$, and the
complete sum $H_7$, which fixes the $\ubar^7$ term of $\eta$ (209 integrands), reaches a relative uncertainty of
$10^{-5}$ in about \costTm{h7} GPU hours, against $5.5\times10^{-5}$ for the same quantity derived from the Phi4tools
series (Sec.~\ref{sec:results}). The new coefficients, from ninth to eleventh order, have relative uncertainties between
\newrelmin\ and \newrelmax\ at $N=1$ (Tables~\ref{tab:newcoeff}--\ref{tab:newcoeffnu}), and the coefficient of
$\ubar^8$ in $\eta$ is now \etaEightGain\ times more precise than before (Sec.~\ref{sec:results-new}). The
uncertainties thus hardly grow with the order. Along the per-graph route the $p^2$ derivatives are the least precise:
at the eighth order of Ref.~\cite{Phi4tools2024}, $5\times10^{-4}$ for $\eta$ against $4\times10^{-6}$ for $\beta$
(Table~\ref{tab:sota}). A uniform precision ensures that the number of known orders, not their statistical errors,
limits the resummation: at the present orders the statistical errors are already negligible there
(Sec.~\ref{sec:resummation-errors}), while the precision of the highest coefficients matters for the large-order and
consistency tests and for the orders still to come.

The paper is organized as follows. Section~\ref{sec:methods} describes the series, the graphs, the numerical
method and the resummation; Sec.~\ref{sec:results} gives the validation, the performance, the new coefficients and
the exponents, compared with other determinations; Sec.~\ref{sec:summary} summarizes. Five appendices give the conventions, the graphs, the zero-dimensional sum rules, the reproducibility checks, and
the availability of the data and the code.
